\section{现代回归：状态空间模型}

\subsection{RNN 思想的复兴}

尽管 Transformer 在过去几年统治了深度学习，但 RNN 的核心思想正在以\textbf{状态空间模型}（State Space Models, SSMs）的形式回归。

\begin{figure}[H]
\centering
\includegraphics[width=0.95\textwidth]{slides-images/slide-052.jpg}
\caption{状态空间模型的代表：RWKV 和 Mamba 等模型结合了 RNN 的线性复杂度和 Transformer 的并行训练能力\protect\footnotemark}
\end{figure}
\footnotetext{Slides 第53页。}

状态空间模型相比 Transformer 的关键优势：

\begin{enumerate}
    \item \textbf{无限上下文长度}：像 RNN 一样，没有固定的上下文窗口限制
    \item \textbf{线性时间推理}：计算量与序列长度成线性关系，而非 Transformer 的二次方
    \item \textbf{固定内存占用}：推理时只需维护固定大小的状态，不需要 KV cache 的线性增长
\end{enumerate}

\begin{importantbox}{RNN vs Transformer 的核心权衡}
\begin{itemize}
    \item \textbf{Transformer}：训练可并行（快），推理时 KV cache 随序列增长（慢），$O(n^2)$ 注意力开销
    \item \textbf{RNN/SSM}：训练难以并行（慢），推理时固定状态（快），$O(n)$ 线性开销
    \item 当前研究热点：如何\textbf{兼得两者优势}——训练时的并行性 + 推理时的效率
\end{itemize}
\end{importantbox}

\subsection{本章小结}

以 Mamba 和 RWKV 为代表的状态空间模型正在挑战 Transformer 的统治地位，特别是在长上下文场景中。RNN 的核心概念——循环状态更新、线性时间推理——正以更现代的形式焕发新生。

%% ========== 拓展阅读 ==========

\subsection{拓展阅读}

\begin{itemize}
    \item Andrej Karpathy: \textit{The Unreasonable Effectiveness of Recurrent Neural Networks}, 2015 \url{http://karpathy.github.io/2015/05/21/rnn-effectiveness/}
    \item Hochreiter \& Schmidhuber: \textit{Long Short-Term Memory}, Neural Computation, 1997
    \item Gu \& Dao: \textit{Mamba: Linear-Time Sequence Modeling with Selective State Spaces}, 2023 \url{https://arxiv.org/abs/2312.00752}
    \item Peng et al.: \textit{RWKV: Reinventing RNNs for the Transformer Era}, 2023 \url{https://arxiv.org/abs/2305.13048}
\end{itemize}

%% ========== 总结与延伸 ==========

